problem:  
Let \( M^n \) be a compact, simply connected Riemannian manifold with nonnegative sectional curvature. Suppose the **symmetry rank** of \( M \) (i.e., the rank of its isometry group) satisfies  
\[
\mathrm{symrank}(M) = \left\lfloor \frac{n}{2} \right\rfloor - 1,
\]
that is, one less than the maximal possible value for such manifolds (as established by Grove–Searle).  

**Problem:**  
Classify, up to diffeomorphism, all such manifolds \( M \). In particular:  
- Do new diffeomorphism types arise that are not homotopy equivalent to known examples such as the compact rank-one symmetric spaces (CROSS) and their quotients by free isometric finite group actions?  
- Can every such \( M \) be realized as a torus manifold or as a biquotient endowed with a nonnegatively curved metric?  
- What constraints exist on the topology of \( M \) (e.g., Betti numbers, cohomogeneity, or Euler characteristic) at this “almost maximal” symmetry rank level?  

Why is it a "good" problem:  
This problem precisely targets the next open frontier following the Grove–Searle classification of manifolds with **maximal symmetry rank**, replacing an incorrect numerical notion of “almost isotropy-maximal” with the correct geometric invariant: **symmetry rank**. The maximal case is fully understood and leads only to CROSS spaces, but the “one-lower” case remains open and is expected to reveal new geometric and topological phenomena.  

It is a “good” problem because:  
1. **It stands at a natural boundary of known theory.** It asks what happens just beyond the known classification—where existing methods (orbit space analysis, torus actions, Alexandrov geometry) almost but not quite suffice.  
2. **It unites multiple domains.** Solving it would require deep integration of Riemannian geometry, transformation group theory, and global topology, especially techniques involving torus actions, cohomogeneity, and equivariant classification.  
3. **It is conceptually simple yet profoundly rich.** “Classify nonnegatively curved manifolds with almost maximal symmetry rank” can be stated in one line, but encompasses a major class of difficult open problems.  
4. **It has high potential yield.** Any new examples or rigidity results discovered in this regime would immediately expand our understanding of how curvature and symmetry constrain topology and could influence the search for new nonnegatively curved manifolds.  

This problem thus exemplifies simplicity, depth, and cross-field resonance—core hallmarks of a high-quality, research-level mathematical question.